\section{Secure Search}
\subsection{Proof of Lemma~\ref{lem:4_1}}\label{app:proof41}

\begin{lemma}[4.1]
Consider a Bloom filter with false positive rate $\frac{1}{m}$, where
$m$ is an arbitrary positive integer. Suppose at most $m$ $\sf BF.Check$ operations
are performed in the BF.
%
Then, for any $\delta > 0$, we have:
$$\Pr [ \mbox{\rm \# false positives} \ge 1+\delta] 
  \le {\frac {e^\delta} {(1+\delta)^{(1+\delta)}}}.$$ 
\end{lemma}
\begin{proof}
Let $\alpha_i$ be the $i$th item that is checked through $\sf
BF.Check$. That is, we consider a sequence of $${\sf
BF.Check}(\alpha_1), \ldots, {\sf BF.Check}(\alpha_{m}),$$
where $\alpha_i$ is an arbitrary item. 
Since we wish to upper bound the false positives (i.e., we don't care about true
positives), it suffices to consider the
case that for every $i$, 
$\alpha_i \not \in {\sf BF}$ (i.e, $\alpha_i$ has not been inserted in the
BF) as this maximizes the number of possible false positives.

Let $X_1, \ldots, X_m$ be independent Bernoulli random variables with
$\Pr[X_i = 1] = 1/m$. 
%
Since the BF false positive rate is assumed to be $1/m$, we have
for all $i$, 
$$\Pr[{\sf BF.Check}(\alpha_i)=1] = \Pr[\mbox{query $i$ is a false positive}] \le
1/m.$$
Thus, we can bound the number of false positives
by $\sum_{i=1}^m X_i$.

Now, let $\mu :=
\Exp[\sum X_i] = m \cdot \frac 1 m = 1$.  By applying the Chernoff bound with $\mu =
1$,  we have: 
%
$$ \Pr\left[\sum_{i=1}^m X_i \ge 1 + \delta \right]  
\le {\frac {e^\delta} {(1+\delta)^{(1+\delta)}}} .$$
\end{proof}


\subsection{Proof of Theorem~\ref{thm:ssp-code}}\label{app:proof63}

\begin{theorem}[6.3]%\label{thm:ssp-code}
Given an FHE scheme, and a $(n, s, c, f_p)$-CODE scheme over domain $D$
in the random oracle model, the construction in Algorithm~\ref{alg:ssCODE} yields a $(\ell,
f_p)$-secure search scheme for records in domain $D$ in the random oracle
model, where $\ell = \frac{c(s)\cdot \ell_c}{s}$, $\ell_c$ is the length
of an FHE ciphertext with plaintext space $D$, and $s$ is the number of
matching records.
\end{theorem}

\begin{proof}
We begin by proving that the adversary cannot distinguish
between two different queries.
The adversary chooses a database $x$ and two queries $q^0, q^1$, with
the promise that $s = \sum_{i=1}^n q^0(x_i) = \sum_{i=1}^n q^1(x_i)$.

The entire view of the adversary during the experiment can be
reconstructed efficiently given (1) the encrypted database $\lbra x
\lket$, (2) the encrypted query $\lbra q \lket$, (3) the decrypted value
of $s$.

We note that the CODE scheme may return either more than $s$ values to the client (in case of a false positive) or less than $s$ values (in case decoding fails), but both of these occur with probability at most $\negl(\secparam)$ and thus we can ignore them in the following.

Since the value of $s$ is the same for $q_0$ and $q_1$, the only thing
that changes in the view of the adversary when switching from $b = 0$ to
$b=1$ is the encrypted query $\tilde{q_b}$.  Therefore, the adversary
guesses $b$ with negligible advantage by the IND-CPA security of the FHE
scheme.

The proof that the adversary cannot distinguish between the same query
applied to two different databases follows nearly identically.

\end{proof}
